\paragraph{Accuracy and generalization (RQ2).}
Table~\ref{tab:accuracy} reports validation-selected checkpoints. Muon+AdamW has the highest mean test accuracy (\pmv{89.04}{0.29}\%). All values use three seeds, not repeated evaluations of one model. The final-epoch column checks whether checkpoint selection changes the interpretation.


\paragraph{Convergence (RQ1).}
The first observed 85\% validation crossing occurs at mean epochs 2.33 (SGD-M), 1.33 (Adam), 1.33 (AdamW), 1.33 (Muon+AdamW). AdamW is fastest by the corresponding training-plus-validation time. Figure~\ref{fig:convergence} shows that ranking by early progress need not match ranking by the final checkpoint. Epoch-end training loss is evaluated on the full training subset, avoiding the changing-model average implicit in an online minibatch loss.
\paragraph{Computational efficiency (RQ3).}
SGD-M has the lowest 15-epoch training cost. Muon takes 1.97$\times$ AdamW time (22.04 versus 11.16 seconds). The extra matrix operations and auxiliary optimizer calls are plausible causes of the overhead, but epoch timing alone does not attribute cost to individual kernels. All methods have the same samples per epoch; throughput therefore expresses the same cost on a per-sample scale.


At the five-second readout, validation-best-so-far checkpoints achieve test accuracies of \pmv{86.94}{0.52}\% (SGD-M), \pmv{87.21}{0.56}\% (Adam), \pmv{87.54}{0.12}\% (AdamW), \pmv{87.18}{0.57}\% (Muon+AdamW). AdamW has the highest mean under this specific accounted-time budget. The budget selects different training durations; this is complementary to, rather than interchangeable with, the fixed-epoch comparison. A strict online deadline, data-loading cost, or a larger model could alter the ranking.
\paragraph{Learning-rate sensitivity (RQ4).}
Across the three tested rates, the best-validation range spans 0.46 points (SGD-M), 1.04 points (Adam), 0.22 points (AdamW), 1.28 points (Muon+AdamW). This range measures sensitivity within the specified grids, not over all rates. A single tuning seed can favor a configuration by chance; three equal trials do not guarantee equally comprehensive searches for different algorithms. Muon has an additional fixed auxiliary learning rate whose sensitivity is not explored. Its winning main rate, .005, lies at the lower grid boundary, so a better rate outside the grid remains possible.
\paragraph{Interpretation and limitations.}
Muon changes mean test accuracy relative to AdamW by +0.45 points while requiring 1.97$\times$ the training time. Its approximate orthogonalization changes update geometry, but that mechanism alone does not imply a gain on this dataset. Higher training accuracy is not sufficient evidence of better generalization: the train--test gaps in Table~\ref{tab:accuracy} must be read with test accuracy, not minimized in isolation. The gaps compare accuracies at the same checkpoint, not training accuracy averaged during optimization.
The largest confounds are parameter routing (only 39.1\% of parameters use Muon), fixed regularization strength, a short horizon, one dataset, and the small number of seeds. A hybrid-control experiment with identical auxiliary AdamW for every method would help isolate the hidden-matrix optimizer. The current experiment supports configuration-level conclusions only. Muon reaches 94.03\% mean training accuracy with a 4.99-point train--test gap, compared with AdamW at 92.68\% and 4.09 points. The additional fitting transfers only partly to test accuracy. Adam and AdamW cross 85\% validation at nearly the same measured time (0.99 versus 0.98 seconds), far smaller than their seed variation; this is not a meaningful speed separation. Sample SD reflects seed variation on one split and GPU; it excludes uncertainty over data splits, hardware, implementation choices, and hyperparameter search. No statistical significance or general superiority is claimed.
